\documentclass[11pt]{article}
\usepackage[margin=0.95in]{geometry}
\usepackage{amsmath,amssymb,booktabs,array}
\usepackage[T1]{fontenc}
\usepackage[hidelinks]{hyperref}
\setlength{\parindent}{0pt}
\setlength{\parskip}{5pt}
\setlength{\emergencystretch}{2em}
\newcommand{\KL}{\operatorname{KL}}
\newcommand{\tr}{\operatorname{tr}}
\newcommand{\norm}[1]{\left\lVert#1\right\rVert}
\newcommand{\pos}[1]{\left[#1\right]_+}
\newcommand{\RR}{\mathbb R}
\title{JLZ v5: Joint Local Targets Through the Actual Writer}
\author{ODE-edit method specification}
\date{October 2, 2026}
\begin{document}
\maketitle

\section{Scope and claim}
All eligible write layers are optimized jointly. Concentration in one
layer is permitted, as is distribution across layers. Current editing
sentences, targets, native readouts and reductions remain those of the
baseline for the same model and benchmark. Preservation reuses only
previously observed native sentences. No external reference corpus,
evaluation paraphrases or neighborhoods, or future editing requests enter
the optimization.

The change from v4 is to optimize a single physical weight model and
commit precisely that model. At each batch entry, a full-context writer
basis is constructed once and then frozen. This is the definition of v5,
not an exact current-key triangular writer followed by an approximation.
The document specifies a method, not a validated implementation or
measured locality improvement.

\section{Variables and fixed writer}
Let $W_t$ be the batch-entry model, $B_t$ the actual logical batch size,
and $\mathcal L$ the complete ordered eligible layer set. For each layer,
\[
D_l\in\RR^{d_{\mathrm{out},l}\times B_t},\qquad
a_{lr}=\norm{h^0_{lr}}>0,\qquad v_{lr}=D_{lr}/a_{lr}.
\]
$D_{lr}$ is a nominal local-target residual. It is also a low-rank write
coordinate; it is not asserted to equal the realized activation change.
With native key-context weights $\sum_c\alpha_{rc}=1$, define
\begin{align}
\bar k_{lr}&=\sum_c\alpha_{rc}k^0_{lrc},\\
A_{tl}&=\lambda_C C_{0l}+H_{tl},\\
M_{tl}&=A_{tl}+\sum_{r,c}\alpha_{rc}k^0_{lrc}(k^0_{lrc})^\top,\\
P_{tl}&=\operatorname{solve}(M_{tl},\bar K_l),\qquad U_l=D_lP_{tl}^\top.
\end{align}
This parameterizes the solution of the entry problem
\[
\min_U\frac12\sum_{r,c}\alpha_{rc}\norm{Uk^0_{lrc}-D_{lr}}^2
+\frac12\tr(UA_{tl}U^\top).
\]
The covariance of individual contexts is retained. The number of columns
of $D$ and $P$ remains $B_t$. A context-dual Woodbury solve or an equivalent
primal solve is allowed, without dropping off-diagonal terms.

For $C=[\sqrt{\alpha_{rc}}k^0_{lrc}]$ and incidence
$R_{(r,c),r}=\sqrt{\alpha_{rc}}$, one has $CR=\bar K$ and
\[
T=A^{-1}C,\quad S=I+C^\top T,\quad P=T S^{-1}R.
\]
All inverses denote solves. Geometry uses FP64 and requires positive
definiteness; silent jitter or removal of layers is not allowed.

\section{Physical graph and exact commit contract}
The reference implementation materializes, once per candidate,
\[
U_l^{64}=D_l^{64}(P_l^{64})^\top,\qquad
W_{\mathrm{eff},l}^{32}=W_{t,l}^{32}+\operatorname{cast}_{32}(U_l^{64}).
\]
Each modified linear operation runs once with this weight. Every native
and replay row sees every request column of the update, at every token.
The final accepted materialized tensors are copied at commit, with no
new solve, gain inversion, or post-fit rescaling.

In the physical graph, the own-layer local action and realized output are
\[
u_{lrc}(D)=D_lP_l^\top k_{lrc}(D_{<l}),\qquad
z^{\mathrm{real}}_{lrc}=h^{\mathrm{pre\ own}}_{lrc}(D_{<l})+u_{lrc}(D).
\]
Lower-layer propagation, cross-request interference, and effects on
non-subject tokens therefore enter the task gradient. $P$ is a constant
by definition; there is no missing $dP/dD$ term. The nominal target
$h^0+D$ need not equal $z^{\mathrm{real}}$. Native-form regularization of
$D$ is not literal native regularization of the realized hidden change.
Mathematically this is also fixed-right-basis low-rank update optimization.
The local-target interpretation comes from deriving that basis from
native contexts and sequential history. The name ``joint-z'' alone is
not a distinction from this equivalent parameterization or a novelty claim.

\section{Native task and two arms}
Let $L_{N,r}$ be the exact native target-token/context-mean NLL, and let
$L_{K,r}$ use the native KL readout with full-vocabulary direction
current $\Vert$ own entry. Define
\[
\bar J_A(v)=\frac1{B_t}\sum_r\left[
L_{N,r}(W_{\mathrm{eff}})+\lambda_{KL}L_{K,r}(W_{\mathrm{eff}}\Vert W_t)
+\sum_l\frac{\lambda_{\mathrm{norm}}}{a_{lr}}\norm{v_{lr}}\right],
\quad\norm{v_{lr}}\le c.
\]
Arm A uses this objective. Arm B adds
\[
\bar J_B=\bar J_A+\lambda_K R_K+\lambda_E R_E.
\]
The v4 quadratic writer-value penalty is absent. There is no additional
shared L1/L2 hard budget, forced balancing, minimum layer count, or
quality-based layer gate. Both arms retain the same ridge geometry and
native-form norm. This preserves loss ingredients and input semantics,
but changes the optimization graph and solver from native compute-z.

\section{Native-only preservation}
A bounded memory stores previously observed facts. A fixed subset $S$
is sampled once per batch and reused throughout its candidates.
Current fact IDs are excluded to permit requested version changes.

For each retained native KL identity, $\pi^{\mathrm{anchor}}$ is an
immutable pre-edit distribution captured at admission. Equal resident
token/mask/position/readout identities share their earliest resident
teacher. An identity lost by eviction can obtain a new admission teacher
through a newly observed fact; a nonexistent old teacher is not assumed
recoverable. This is not a global $W_0$ teacher.
\[
R_K=\frac1{|S|}\sum_{r\in S}\sum_q\beta_{rq}
\KL\big(p_{W_{\mathrm{eff}}}(\cdot\mid x^{KL}_{rq})
\Vert\pi^{\mathrm{anchor}}_{rq}\big).
\]
For native rewrite context $c$, let $b_{rc}$ be the target-token-mean NLL
at the fact's latest actual commit. With native rewrite weights $w_{rc}$,
\[
R_E=\frac1{|S|}\sum_{r\in S}\sum_c w_{rc}\,
\frac12\pos{\ell_{rc}(W_{\mathrm{eff}})-b_{rc}}^2.
\]
Improvement in past target NLL is unpenalized. This does not preserve the
entire postcommit distribution, and needs only scalar rewrite anchors.
Weak committed edits produce weak anchors; admission is not conditioned
on edit success. Anchors are updated only upon an explicit new version
of that fact, not at every batch entry. Empty-memory terms are zero.
Consequently the first-batch objectives of A and B coincide.

Algorithm R samples first-observed unique fact IDs into a memory of
capacity $M$. Repeated in-memory facts retain their slot and update only
their rewrite inputs, target and NLL anchors. Their original KL input
tuple and teacher remain frozen together. Previously seen out-of-memory
facts are not readmitted. The seen-ID metadata grows with distinct facts;
$M$ bounds payload count, not all metadata. Shared teachers use owner
reference counts, with eviction and admission applied in stream order.
Pending records become visible after commit only. Exclusion uses current
fact IDs; another retained fact can still share a current KL prefix.
Selection depends on stream metadata and a fixed seed, never on scores.

Provisional Llama/CounterFact settings are $M=128$,
$|S|=\min(16,B_t,\text{eligible memory size})$,
$\lambda_K=\lambda_{KL}=0.0625$, and $\lambda_E=1$.
These are explicit unvalidated choices, not universal model constants.

\section{Normalization and proximal optimization}
For request-SUM implementation use
\[
f(v)=\sum_r(L_{N,r}+\lambda_{KL}L_{K,r})
+\eta B_t(\lambda_K R_K+\lambda_E R_E),\qquad
\Omega(v)=\sum_{l,r}w_{lr}\norm{v_{lr}},\quad
w_{lr}=\lambda_{\mathrm{norm}}/a_{lr}.
\]
Here $\eta\in\{0,1\}$. Scaling the whole mean objective by $B_t$ preserves
its minimizers; adding an unscaled reference mean to a native sum would
not preserve the relative loss units.

With a single scalar step $\tau$, $g=\nabla f(v)$, and
$y_{lr}=v_{lr}-\tau g_{lr}$, the exact group-norm-plus-ball prox is
\[
v'_{lr}=\frac{y_{lr}}{\norm{y_{lr}}}
\min\{c,\pos{\norm{y_{lr}}-\tau w_{lr}}\},
\]
with zero output for zero $y$. Every layer gradient is computed, so a
zero group may reactivate. No per-group unit-gradient normalization or
momentum is used. Initially take $\nu=c/4$ and
$\tau=\nu/\max_{lr}\norm{g_{lr}}$; zero gradient at zero returns zero.
The zero candidate is the initial accepted state. Later trial steps are
no larger than the last accepted step and this smooth-displacement bound.
Halve the active step after rejection, retaining the accepted state and
gradient; do not reset the step before the next retry.

Accept a trial using smooth majorization,
\[
f(v')\le f(v)+\langle g,v'-v\rangle+
\frac{\norm{v'-v}^2}{2\tau}+\epsilon_{\mathrm{num}}.
\]
The numerical tolerance is bound by runtime qualification, not PS/NS.
This is objective acceptance, not a locality or allocation gate.
The proposed budget is 25 candidate evaluations, including zero and
rejected trials: candidates 1--24 include backward, candidate 25 is
forward-only, and the last accepted state is returned. Accepted updates
can be fewer than 24. Early stationary exit records actual counts and
uses the proximal mapping or equality of $v$, not merely identical
rounded materialized weights.

\section{History, efficiency, and limits}
After commit, append each occurrence once using actual context keys:
\[
H_{t+1,l}=H_{t,l}+\sum_{r,c}\alpha_{rc}
k^{\mathrm{commit}}_{lrc}(k^{\mathrm{commit}}_{lrc})^\top.
\]
Reuse accepted-forward keys only when the adapter proves causal key-input
identity. Entry keys are not automatically commit keys. Older history
keys can still become stale; this design does not refresh all past data.
Physical replay measures current drift on sampled native inputs only.

Cache only computations before the earliest modified linear operation.
Upper-layer subject-prefix KV caches are invalid. Fixed $P$ permits
independent microbatch graph release with a joint accumulated gradient.
For linear input $X$ and output gradient $G$,
\[
\nabla_D L=G^\top(XP),\qquad \nabla_X L=G W_{\mathrm{eff}}.
\]
Dense weight gradients are unnecessary, but reordered backward arithmetic
requires comparison with the dense reference. Selected-position heads
retain the full vocabulary. Entry preparation, context solves,
materialization, rejected trials, replay, and commit are all timed.

The interface supports variable logical batches, context counts, layer
dimensions, and model/benchmark profiles. This does not imply identical
results after changing batch boundaries or universal hyperparameters.
No performance gain follows solely from fewer optimized coordinates.
Native-only references do not guarantee protection of unseen neighbors,
all earlier facts, or general knowledge relative to $W_0$.

Weights, history, memory, RNG, and pending admissions form one RAM
transaction. Existing no-checkpoint policy is retained. This design
does not modify or resume an ongoing v4 chain.
\end{document}
